\subsection{TOPOLOGICAL MODULES}

DEFINITION 3. Given a topological ring A with an identity element, a topological left A-module is a set E, together with:

\begin{enumerate}
\def\labelenumi{\arabic{enumi})}
\item
  a left A-module structure;
\item
  a topology compatible with the additive group structure of E, and satisfying the following axiom:
\end{enumerate}

(MT) The mapping \((\lambda, x) \to \lambda x\) of \(\mathbf{A} \times \mathbf{E}\) into \(\mathbf{E}\) is continuous.

We define similarly the notion of a topological right A-module; since every right A-module can be considered as a left A°-module, where A° is the opposite ring of A, and since the topology of A is compatible with the ring structure of A°, there is no need to distinguish between topological right A-modules and topological left A°-modules.

Examples. *1) A topological vector space over R (resp. C) is a topological R (resp. C) module.

\begin{enumerate}
\def\labelenumi{\arabic{enumi})}
\setcounter{enumi}{1}
\tightlist
\item
  Let A be a ring and let B be a filter base on A consisting of two-sided ideals of A; let E be a left A-module. If we give A the topology (compatible with its ring structure) for which B is a fundamental system of neighbourhoods of o (no. 3, Example 3), and E the topology (compatible with its additive group structure) in which the sets aE, as a runs through B, form a fundamental system of neighbourhoods of o (§ 1, no. 2, Example), it is immediately verified that E is a topological A-module.
\end{enumerate}

Remark. Given a topological ring A, consider, on a left A-module E, a topology compatible with the additive group structure of E. By virtue of the identity

\[
\lambda x - \lambda_ {0} x _ {0} = (\lambda - \lambda_ {0}) x _ {0} + \lambda_ {0} (x - x _ {0}) + (\lambda - \lambda_ {0}) (x - x _ {0})
\]

the axiom (MT) is equivalent to the conjunction of the following three axioms:

\((\mathbf{MT}_{\mathbf{I}}^{\prime})\) . For each \(x_0\in \mathbf{E}\) , the mapping \(\lambda \to \lambda x_0\) is continuous at the point \(\lambda = 0\)

\((\mathbf{MT}_{\mathrm{II}}^{\prime})\) . For each \(\lambda_0\in \mathbf{A}\) , the mapping \(x\to \lambda_0x\) is continuous at the point \(x = 0\)

\((\mathbf{MT}_{\mathrm{III}}^{\prime})\) . The mapping \((\lambda ,x)\to \lambda x\) is continuous at the point (o, o).

We deduce from this a necessary and sufficient set of conditions that the filter \(\mathfrak{B}\) of neighbourhoods of o in an A-module E must satisfy in order to define a topology on E compatible with its module structure; \(\mathfrak{B}\) must satisfy the axioms (GA \(_{\mathrm{I}}\) ) and (GA \(_{\mathrm{II}}\) ) of § 1, no. 2, and in addition must satisfy the following three axioms:

\((\mathbf{MV}_{\mathrm{I}})\). For each \(x_0 \in E\) and V \(\in \mathfrak{B}\) , there is a neighbourhood S of o in A such that S. \(x_0 \subset V\) .

\((\mathbf{MV}_{\mathrm{II}})\) . For each \(\lambda_0 \in \mathbf{A}\) and \(V \in \mathfrak{B}\) , there exists \(W \in \mathfrak{B}\) such that \(\lambda_0 W \subset V\) . \((\mathbf{MV}_{\mathrm{III}})\) . For each \(V \in \mathfrak{B}\) there exists \(U \in \mathfrak{B}\) and a neighbourhood \(T\) of \(o\) in \(A\) such that \(T. U \subset V\) .

Every commutative topological group is a topological \(\mathbf{Z}\) -module when the ring \(\mathbf{Z}\) is given the discrete topology.

If M is a submodule of a topological A-module E, it is clear that the topology induced on M by the topology of E is compatible with the module structure of M. Moreover, on the quotient A-module E/M, the topology which is the quotient by M of the topology of E is compatible with the A-module structure. To see this it is enough to show that the mapping \((\lambda, \dot{x}) \to \lambda \dot{x}\) of \(\mathbf{A} \times (\mathbf{E} / \mathbf{M})\) onto \(\mathbf{E} / \mathbf{M}\) is continuous (where \(x \to \dot{x}\) denotes the canonical mapping of \(\mathbf{E}\) onto \(\mathbf{E} / \mathbf{M}\) ). Now, since we may identify the additive topological groups \(\mathbf{A} \times (\mathbf{E} / \mathbf{M})\) and \((\mathbf{A} \times \mathbf{E}) / (\{\mathbf{o}\} \times \mathbf{M})\) (§ 2, no. 9, Corollary to Proposition 26), it is enough to show that the mapping \((\lambda, x) \to \lambda \dot{x}\) of \(\mathbf{A} \times \mathbf{E}\) into \(\mathbf{E} / \mathbf{M}\) is continuous; and this is immediate, since the mapping in question is the composition of \(x \to \dot{x}\) and \((\lambda, x) \to \lambda x\) .

Let \((\mathbf{E}_i)_{i\in \mathbf{I}}\) be an arbitrary family of topological A-modules, and let \(\mathbf{E} = \prod_{i\in \mathbf{I}}\mathbf{E}_i\) be the A-module which is the product of the \(\mathbf{E}_i\) . Then the

product topology on E is compatible with the A-module structure of E. To prove this it is enough to show that the mapping \((\lambda, x) \to (\lambda \cdot \text{pr}_i x)_{i \in I}\) of A × E into E is continuous, or (by Proposition 1 of Chapter I, § 4, no. 1) that for each index \(i \in I\) the mapping \((\lambda, x) \to \lambda \cdot \text{pr}_i x\) is a continuous mapping of A × E into \(E_i\) ; but this mapping is the composition of \((\lambda, x_i) \to \lambda x_i\) and \((\lambda, x) \to (\lambda, \text{pr}_i x)\) , both of which are continuous.

Let A be a Hausdorff topological ring and E a Hausdorff topological A-module. Let \(\hat{\mathbf{E}}\) be the additive group which is the completion of the commutative topological group E (§ 3, no. 5, Theorem 2). The Z-bilinear mapping \((\lambda, x) \to \lambda x\) of the product \(\mathbf{A} \times \mathbf{E}\) of the additive groups A, E into the additive group E can be extended by continuity to a Z-bilinear mapping of \(\hat{\mathbf{A}} \times \hat{\mathbf{E}}\) into \(\hat{\mathbf{E}}\) (no. 5, Theorem 1), and this mapping we continue to denote by \((\lambda, x) \to \lambda x\) . By virtue of the principle of extension of identities, we have \(\lambda(\mu x) = (\lambda \mu)x\) for \(\lambda \in \hat{\mathbf{A}}\) , \(\mu \in \hat{\mathbf{A}}\) and \(x \in \hat{\mathbf{E}}\) , and \(1.x = x\) for all \(x \in \hat{\mathbf{E}}\) ; the external law \((\lambda, x) \to \lambda x\) therefore defines an \(\hat{\mathbf{A}}\) -module structure on \(\hat{\mathbf{E}}\) compatible with its topology. The topological \(\hat{\mathbf{A}}\) -module \(\hat{\mathbf{E}}\) thus defined is called the completion of the topological A-module E.

Let E be a topological module over a topological ring A, where neither A nor E is necessarily Hausdorff. Let N (resp. F) be the closure of \(\{o\}\) in A (resp. E). N is a two-sided ideal of A (no. 4, Proposition 5) and F is a sub-A-module of E (no. 1, Proposition 1); furthermore, by continuity we have \(\lambda x \in F\) whenever \(\lambda \in N\) or \(x \in F\) . We can therefore define, by passing to the quotients, a mapping \((\dot{\lambda}, \dot{x}) \to \dot{\lambda} \dot{x}\) of \((A/N) \times (E/F)\) into E/F; it is easily verified (by use of the Corollary to Proposition 26 of § 2, no. 9) that this mapping is continuous, and therefore defines a structure of a topological \((A/N)\) -module on E/F. If we put B = A/N and L = E/F, then the B-module L is called the Hausdorff module associated with E; its completion \(\hat{L}\) is a topological module over the Hausdorff completion \(\hat{A}\) (equal by definition to \(\hat{B}\) ) of A (no. 5), and this module \(\hat{L}\) is called the Hausdorff completion of E and is denoted

by \(\hat{\mathbf{E}}\) . We see as in § 3, no. 4, Proposition 8 that every continuous homomorphism \(u: \mathbf{E} \to \mathbf{G}\) of \(\mathbf{E}\) into a complete Hausdorff \(\hat{\mathbf{A}}\) -module \(\mathbf{G}\) factorizes uniquely into \(u = v \circ \varphi\) , where \(v\) is a continuous homomorphism of \(\hat{\mathbf{E}}\) into \(\mathbf{G}\) and \(\varphi\) is the canonical mapping of \(\mathbf{E}\) into \(\hat{\mathbf{E}}\) . We conclude that if \(\mathbf{E}, \mathbf{E}'\) are two topological A-modules and \(u: \mathbf{E} \to \mathbf{E}'\) is a continuous homomorphism, then there is a unique continuous homomorphism \(\hat{u}: \hat{\mathbf{E}} \to \hat{\mathbf{E}}'\) such that the diagram

\[
\begin{array}{c c c} \mathbf {E} & \stackrel {{u}} {{\longrightarrow}} & \mathbf {E} ^ {\prime} \\ \varphi \Big \downarrow & & \Big \downarrow \varphi^ {\prime} \\ \hat {\mathbf {E}} & \stackrel {{\hat {u}}} {{\longrightarrow}} & \hat {\mathbf {E}} ^ {\prime} \end{array}
\]

is commutative, \(\varphi\) and \(\varphi'\) being the canonical mappings.

The definitions and results of this section apply equally well to pseudo-modules over an arbitrary topological ring, by deleting all mention of the identity element of the ring.

